\documentclass[12pt]{article}
%friendlyname:Niyagh;Plants
\usepackage{covington,graphicx,tipa,makeidx,fontspec}
\setmainfont[Mapping=tex-text]{Charis SIL}

\begin{document}

\title{Talking about plants}
\author{Belle Deacon}
\date{}
\maketitle

Recording by Jim Kari.  John Deacon also present.

Audio:  portion of anl0770.wav. There is a loud noise in the middle of the recording which has been edited out here.

Length of edited recording:  7:43.

Translation by Edna Deacon and Lucy Hamilton.

Transcription by Sharon Hargus.

\bigskip
This is a set of short recordings about certain plants.  The recording
clicks off and on again between topics.

\section{Berries}

\gll Iy getiy viq'aditr'et'a \em{that} niɬanht'asr.
it really {we love it} {} blackberry
\glt `We really love those blackberries.'
\glend

cause

\gll nenhtl'it chenh
{low-bush cranberry} too
\glt `low-bush cranberries too'
\glend

\gll Viqa xeydantr'ididix     gochexi
{off it} {we live in winter} ?
\glt `We live on them all winter.'
\glend
%ED:  gochexi pronounciation recognized, but not translated

and uh

\gll Sanh tth'ixinedr xits'i xodhiɬ tux cheche dandhu'on che yitr'inehayh, iy zro iy axa--- 
summer middle  {to it} {time is passing} when again salmonberry too {we pick them} it only it {with}
\glt `When it's getting to be the middle of summer we pick salmonberries too, only that.'
\glend

\gll Yitotl'ogh tux che xisrghed yiɬ, axaxiɬdik tr'onihay che.
{after that} when again rosehips too then {high-bush cranberry} again
\glt `After that there are rosehips and high-bush cranberries.'\footnote{a.k.a. ginathdloy}
\glend

\gll Yitots'in' {gag veda} yiɬ.
then       {black currant} too
\glt `Then there are black currants too.'
\glend

\gll Axaxiɬdi nondiniy chenh.
then     {red currant} too
\glt `Red currants too.'
\glend

\gll Axaxiɬdik {dandiggiy, dandiggiy} che, sanh tux dandiggiy yitr'eheyh.
then      raspberry          too  summer in raspberry {we eat them}
\glt `In the summer we eat raspberries too.'
\glend

\gll Dandiggiy getiy ngizrenh.
salmonberry very {it's good}
\glt `Raspberries are awfully good.'
\glend

\gll Ixigide iy go, {go                  daɬts'in}               dinaqay vixitinalyay sanh tux.
{around here} it here {this side (of Alaska)} {our village} {those which are around} summer in
\glt `Those are the ones (berries) that grow around our village in the summer.'
\glend

JD:  that's correct

BD:  'n

\gll Iy, niɬanht'asr iy, edi xeydi 'n edi sanh cheche tso yedathdlo, nenhtl'it yiɬ.
it  blackberry it all winter {} all summer again cache {keep them in} cranberry too
\glt `All winter and all summer you can keep blackberries in the cache with cranberries.'
\glend 

it's good all win--

\gll Sanh tux xiyiɬ vixiyiɬ ngizrenh. 
summer in then {} {it's good} 
\glt `They're good in the summer.'
\glend

\gll {Gits'andithdhigi ts'in'}.
{{it doesn't spoil}}
\glt `They don't spoil.'
\glend

\gll Agidet.
{{that's all}}
\glt `That's all'
\glend

\section{More on berries and other plants}

BD:  like this?

JK:  mhm

BD:  uh

\gll Nenhtl'it \em{is} {dinaye dixineyh} tux, {xigi'imo ghined,} nenhtl'it.
cranberry {} {they get diarrhea}  when {their medicine} cranberry
\glt `Cranberries are good medicine for diarrhea.'\footnote{cooked berries; cook hard for 30 mins. with 1 T. sugar; 2-3x as much water as berries; bit of cornstarch to consistency of thin gravy)}
\glend

\gll Axaxiɬdi {gag veda} cheche, {gag veda} ivatr tux, getiy                {quth di- dina'ileyh} tuxdi xone ngizrenh iy.
then    {black currant} too {black currant} {it simmers} when very {we have cough} when for {it's good} it
\glt `Then stewed black currants are good when we have a bad cough.'\footnote{make pudding out of it?}
\glend

(phone rings)

%(loud noise)

BD:  I'll push this one?

JK:  that's correct

BD:   yeah. 

\gll Xileggi tux che xoɬtthil xintthi'itthiyh tux iy che ivatr tux getiy ngizrenh.
spring in       again rhubarb  {it sticks its head up} when it too {it simmers} when very {it's good}
\glt `In the spring when rhubarb comes out it too is good boiled.'\footnote{ED adds dried apples to it; makes pudding with cornstarch}
\glend

and uh

\gll Ilivi vadr yiɬ, iliv che tr'iɬvatr che. Iy che ngizrenh.
nettle boiled too nettle too {we boil} too it too {it's good}
\glt `Boiled nettles too, we boil them too. That's good too.'
\glend

\gll Agide edi.
{that's all} {all over}
\glt `There's no more.'
\glend

\section{Nettles}

\gll Xeyts'in' tux, iliv   diditritr tux           dughontr'iɬch'iɬ tux.
fall      when nettle {it turns to wood} when {we peel them} when
\glt `In the fall, when nettles becomes woody, we peel them.'
\glend
%ED  ditritr tux

\gll {Yuxudz yeg} tth'ax xiq'i vuqogg yiɬ xelanh iy.
actually sinew like {its surface} too {it is} it
\glt `Its bark is just like sinew.'
\glend

\gll Getiy ngizrenh hiy n'iy.
very {it's good} it it
\glt `It's very good.'
\glend

\gll Yitongo sanh tux getiy {vilitr'aghdithnog ts'i} dinalo', getiy vits'i imo delanhdi xo'in.
meanwhile summer during very {we can't touch it}   {our hand} very  {from it} sickness {there is} because
\glt `In the summer we can't touch it, because it really hurts.'\footnote{it penetrates our skin}
\glend

\gll Agidet.
{{that's all}}
\glt `That's all.'
\glend

\section{More on nettles}

\gll Iliv di'aɬch'iɬ tux getiy vitth'agh \em{just} yeg tth'ax xiq'i.
nettle {it's peeled}   when very {it's sinew} {} {around there}    sinew like
\glt `When nettles are peeled they're just like sinew.'
\glend

\gll Axaxiɬdi iy yiggiy gaɬtim yiɬ axaxiɬdik ni'it'id yiɬ niɬtoxiynitltl'etl tux noghniy gho dughonxiynathtl'enh.
then     it that   {inner willow bark} with then {peeled} with {they braid them together} when rabbit for {they made snares}
\glt `They braided them with peeled inner willow bark and made rabbit snares.'
\glend

\gll Getiy ditth'aq.
very {it's strong}
\glt `It's really strong.'
\glend

\gll Agidet.
{{that's all}}
\glt `That's all.'
\glend

\section{Juniper}

\gll Axaxiɬdik {ttha il} che getiy {imo ghined} ngizrenh.
then       juniper  too very  medicine     {it's good}
\glt `Juniper is good medicine too.'
\glend

\gll {Ttha il} iy yuxudz vigaga' iy yiɬ xelanh iy gag niɬyagh xiq'i.
juniper   it just   {its berry} it too {there is} it berry blueberry like
\glt `Juniper berries look like blueberries.'
\glend

\gll Iy xiɬdi iy {getiy gitr'et'a} tux  giɬiggi dinadhatr'inu'oyh ts'i then(?) tr'u'oɬ vite' tr'italnax ts'in'.
it then  it {we're very sick} when one     {we pop one in the mouth}  and {} {we chew} {its liquid} {we'll swallow} 
\glt `When we're awfully sick we just put one in our mouth and chew and swallow the liquid.'
\glend

\gll {Imo ghined} getiy ngitl'itth.
medicine   very {it's strong}
\glt `It's very powerful medicine.'
\glend

\gll Ngizrenh, imo tr'elanhdi xonet.
{it's good} sick {when we are} for
\glt `It's good for our sickness.'
\glend

\gll Agidet.
{{that's all}}
\glt `That's all.'
\glend

\section{Pitch}
 
\gll Dzax, dinalo', dinaqa' tr'iɬquyh tux xiyiɬ, dzax yit xidoghq'atr'iɬtth'ix.
pitch {our hand} {our foot} {we chop} when then pitch there {we smear it}
\glt `When we chop our hands or our feet, we smear pitch on it.'
\glend

\gll Axaxiɬdi iy yiɬ vagaghditl'eyh tux {vidoq'agithogi ts'in'.}
then     it with {wraps it }          when {it doesn't get infected}
\glt `Then it's wraps it like gauze and it doesn't get infected.'
\glend

\gll \em 'Cause \em dzax getiy dinaqogg a xone ngizrenh.
{}     pitch very {our skin} {} for {it's good}
\glt `Pitch is good for (cuts on) our skin.'
\glend

\gll Axaxiɬdi dzax che dixiɬ'anh vedoy xiyaxa iy ntodzaq hiy.
then     pitch too {they do} boat {they with it} it {will paint on} it
\glt `That's the kind we pick for our boat, which is painted on it.'
\glend

\gll Getiy {xilongh ts'i} vaxa dixit'anh go dzax.
very {many} {with it} {they do} this pitch
\glt `They use it for lots of things.'
\glend

\gll Agide.
{{that's all}}
\glt `That's all.'
\glend

\section{Cottonwood bark}

\gll Yidong tux gilotr'e               dughonxiɬquyh. 
{in the past} when {peeled cottonwood bark} {they chop off}
\glt `In the past they would chop off cottonwood bark.'
\glend

because 

\gll Xiyaxa yix--  xiyaxa {sanh yix} xiyxeghoyh.
{they with it} {} {they with it} {smokehouse} {they make them}
\glt `They make smokehouses out of it.'
\glend

\gll Ngiyixi yix xiyiɬ xiq'i xiydinelo.
floor house then on {they put it}
\glt `They put it down on the ground, as flooring.'
\glend

\gll {Vixidithiq'izr ts'i}    q'idogh,      {vixidithiq'izr ts'in'.}
{it doesn't wear out}  {not right away} {it doesn't wear out} 
\glt `It doesn't wear out right away.'
\glend

\gll {Edi sanhdi}  {vaxa di-- dixit'anh} {xiyqogg nixitset,} go gilotr'e. 
{all summer}  {they use it}         {they sweep it off} this {peeled cottonwood bark}
\glt `All summer long they use it, sweeping it off.'
\glend

\gll Axaxiɬdik ngidosin, ngidosin che {xiyaxa xiɬq'iyh}.
then     roof roof too {they cover with it}
\glt `They make roofing with it too.'
\glend

\gll {Yuxudz yeg} voghyigg te {dithtligi ts'in'}. 
actually           {under it} water {it doesn't drip}
\glt `Water doesn't drip under it.'
\glend

\gll {Getiy, getiy} ngizrenh ts'i dinaɬt'adi sre' nixidodhil xiq'idathdlo, hingo xeɬedz xiq'idi'alyayh tux.  
very {it's good} {} {how many} {I don't know} {passing time} {they keep on top} while well  {they put it on top} when
\glt `I don't know how many years they could keep it on the roof, when it's put on correctly.'\footnote{Shelters were also lined with the sticky inner side, which would trap mosquitos.---ED}
\glend

\gll Agide.
{{that's all}}
\glt `That's all.'
\glend

\section{Spruce}

\gll Ts'ivi'il te' che {imo delanh hinh} xiyetonxidalyax ts'ixiyiɬ.
{spruce branch} liquid too {sick people} {they take a bath in it} then
\glt `Sick people bathe in the liquid from spruce branches.'
\glend

\gll Yitots'i xiyiɬ ts'ixiyiɬ tigitth'ok axa xiyiɬ {xiyits'i dhagiditeyh.}
then then then teaspoon with then {they put it in mouth}
\glt `They even can drink it with a teaspoon.'
\glend

\gll {Ndadz sre'} getiy {imo ghined} tiyh xaɬne' ts'in'.
{I don't know what} very medicine powerful {they said}
\glt `They say it's some kind of really powerful medicine.'
\glend

\gll Agide dixiyiɬ'anh ts'in'.
{that's all} {they use it}
\glt `That's how they used it.'
\glend

\section{Birch}

\gll Go q'eyh xiɬdik getiy  q'eyh tr'itiɬch'iɬ tux tetth'ok yiɬ yitr'eghoyh.
this birch then very birch {we peel} when basket too {we make them}
\glt `When we peel birch we make baskets with it.'
\glend

\gll Axaxiɬdi yiggi q'eyh dixinaɬt'adi sre'           nixidodhil xiyiɬ q'eylgguddh         che {vogh dontr'iɬch'iɬ} {viq'idz gidit'odh} lay.                 
then    that birch {how many times} {I don't know} {years} then  {second growth birch bark} too {we peel them}  {cutting board} material
\glt `How many times, I don't know how many years later (you go back to the same tree) and peel second growth birch bark for (fish) cutting boards.'
\glend

\gll {Xilongh ts'i} {vaxa dixit'anh} go q'eyh.
many            {they use it} this birch
\glt `This birch is used for many things.'
\glend

\gll Go che vedoy yiɬ yixigheghoyh hiy.
this too boat too {they used to make them}
\glt `They use to make canoes with them.'
\glend

\gll Yitots'i go q'eyh cheche iy xutl    yixighegho-- yixigheghoyh ts'i tigitth'og yiɬ.
then     this birch also that sled  {} {they used to make them} and spoon too
\glt `They used to make sleds and spoons from birch too.'
\glend

\gll Q'idoghtux che tritr tigitth'og, tritr tth'ok yiɬ yixiygheghoyh.
sometimes  also wooden spoon wooden bowl {} {they used to make them}
\glt `Sometimes they made wooden spoons and bowls too.'
\glend

\gll Getiy {xilongh ts'i} ghonet dingit'a hiy go q'eyh.
very  many           for   {it is} it this birch
\glt `This birch is used for many things.'
\glend

\gll Agide.
{{that's all}}
\glt `That's all.'
\glend

JK:  q'ingitl (`sap'). qingitl?

BD:  not \em q'ingitl \em 
 
\section{Birch sap}

\gll Q'eyh getiy vivi- q'ingidili-  che xiyghon'iɬquyh         di'iɬt'ayh tux        xiyghoyixi gathdlo.
birch very {} {} also {they chop through} {it's running} when {they under it} {keep something}
\glt `They chop through birch when (sap) is running and keep something under it.'
\glend

\gll Iy xiɬdi iy getiy xiyiɬvatr tux getiy liniyh.
it then  it very {they cook it} when very {it's tasty}
\glt `It tastes good when they cook it on the stove.'
\glend

\gll Liq'ets ts'i  xiyiɬ yixiyiɬvatr.
{it's sweet} then then {they boil them}
\glt `It gets sweet when they cook it.'
\glend

\gll {Anhtr'a nyeg} shaxil q'idingit'a.
{as if}  sugar {it's like}
\glt `It's like sugar.'
\glend

\gll Agide.
{{that's all}}
\glt `That's all.'
\glend

\section{Willow and alder}

\gll Tr'itl, q'eylgesr \em and {tixet'an q'uy} \em and q'isr ixigide vixitinalyay sanh tux.
willow  {willow tip} {} {unidentified willow} {} alder {around here} {the ones around here} summer during
\glt `Willows and alder grow around here during the summer.'
\glend
%BD says vinxitinalyay

\gll Agide.
{{that's all}}
\glt `That's all.'
\glend

JD:  that's it

BD:  no, I couldn't

\section{Inner willow bark}

\gll (Gaɬ)tim ɬegg gitolixdi xits'i iɬt'et tr'anqaye  yigixiniɬtl'etl.
{inner willow bark} fish {will run} towards continually {old women} {they braid it}
\glt `The old women are always braiding inner willow bark when it's time for the fish to run.'
\glend

\gll \em{Cause} iy axa dit'anh iy axa ni'iniyh ts'i.
{} it with {people do} it with {people work with}
\glt `Because people work with it, do things with it.'
\glend

\gll Xiyaxa xiɬdi dughonxiɬch'iɬ ts'in'.
{they with it} then {they tie them together} {}
\glt `They tie things together with it.'
\glend

\gll Xilonghts'i vaxa dixit'anh go, tr'itl gigaɬtim.
{many times} {with it} {they do} this willow {inner willow bark}
\glt `They do lots of things with this inner willow bark.'
\glend

\gll Agide. Idits'in'.
{that's all} {forever}
\glt `That's it. No more.'
\glend

\section{Wild celery}

\gll Dichiyedim iy vits'ixilu'o hiy n'iy deg engithidong xits'i vits'ixilu'o, dichiyedim.
{wild celery} it {it's honored} it {} {around here} {in the past} from {it's honored} {wild celery}
\glt `Wild celery has been an honored plant around here since long ago.'
\glend

\gll \em because \em digathtonh tux {xiyaxa dit'anh}.
{}        {lucky dance is held} when {they use it}
\glt `Because when they had a lucky dance they used it.'
\glend

\gll Dichiyedim yan' q'ixidala ts'i yuxudz axa xiytitl-'o go dichiyedim.
{wild celery} only {they carry around} and all with {they praised it} this {wild celery}
\glt `Only wild celery was carried and praised.'
\glend

\gll Agide. 
{{that's all}}
\glt `That's it. 
\glend

\end{document}